
\documentclass{article}
\usepackage{colortbl}
\usepackage{makecell}
\usepackage{multirow}
\usepackage{supertabular}

\begin{document}

\newcounter{utterance}

\centering \large Interaction Transcript for game `codenames', experiment `risk\_high', episode 6 with Qwen3.6{-}35B{-}A3B{-}FP8{-}without{-}reasoning.
\vspace{24pt}

{ \footnotesize  \setcounter{utterance}{1}
\setlength{\tabcolsep}{0pt}
\begin{supertabular}{c@{$\;$}|p{.15\linewidth}@{}p{.15\linewidth}p{.15\linewidth}p{.15\linewidth}p{.15\linewidth}p{.15\linewidth}}
    \# & $\;$A & \multicolumn{4}{c}{Game Master} & $\;\:$B\\
    \hline

    \theutterance \stepcounter{utterance}  
    & & & \multicolumn{2}{p{0.3\linewidth}}{
        \cellcolor[rgb]{0.9,0.9,0.9}{
            \makecell[{{p{\linewidth}}}]{
                \texttt{\tiny{[GM$|$GM]}}
                {'hidden': {'team': ['medus', 'lochnesas', 'kasykla', 'arklys', 'morka', 'bakstelėkite', 'burna', 'laikas', 'baras'], 'innocent': ['jupiteris', 'erelis', 'šepetys'], 'opponent': ['styga', 'pyragas', 'amerika', 'mėta', 'siena', 'siela', 'degtukas', 'blokas'], 'assassin': ['alyvuogė', 'lokys', 'kengūra', 'slaugytoja', 'maskva']}, 'revealed': {'team': {'team': [], 'innocent': [], 'opponent': [], 'assassin': []}, 'opponent': {'team': [], 'innocent': [], 'opponent': [], 'assassin': []}}}
            }
        }
    }
    & & \\ \\

    \theutterance \stepcounter{utterance}  
    & & \multicolumn{4}{p{0.6\linewidth}}{
        \cellcolor[rgb]{0.9,0.9,0.9}{
            \makecell[{{p{\linewidth}}}]{
                \texttt{\tiny{[P1$\langle$GM]}}
                \texttt{Pateikite vieną vieno žodžio užuominą, susijusią su kai kuriais jūsų komandos žodžiais iš toliau pateikto sąrašo. Galite pasirinkti, su kiek ir kokiais žodžiais bus susijusi jūsų užuomina.} \\
\texttt{Tikslas {-} kad jūsų partneris, remdamasis jūsų pateiktomis užuominomis, atspėtų visus jūsų komandos žodžius, kol kita komanda neatskleidė visų savo žodžių. Jei jūsų partneris atspėja "priešininko" arba "atitraukiančio dėmesį" žodį, jo spėjimo eilė baigiasi ir priešininkų komanda gali atlikti savo ėjimą. Jei jūsų partneris atspėja "žudantį" žodį, žaidimas baigiasi ir jūsų komanda iš karto pralaimi.} \\
\texttt{Kuo daugiau žodžių, su kuriais susijusi jūsų užuomina, tuo greičiau galite laimėti žaidimą, tačiau partneriui gali būti sunkiau atspėti asociacijas, todėl išmintingai pasirinkite užuominą. Užuominos žodis turi būti semantiškai susijęs su tiksliniais žodžiais, jis negali būti vienas iš sąrašuose esančių žodžių arba sudaryti jų dalių.} \\
\\ 
\texttt{Visada pateikite vieno žodžio užuominą ir kableliais atskirtą susijusių tikslinių žodžių sąrašą toliau nurodyta forma ir atsakymus pateikite kuo trumpesnius, niekada nepateikite jokio kito teksto, nei reikalaujama šioje formoje:} \\
\\ 
\texttt{UŽUOMINA: <WORD>} \\
\texttt{TIKSLAI: <TARGETS>} \\
\\ 
\texttt{Jūsų komandos žodžiai yra: medus, lochnesas, kasykla, arklys, morka, bakstelėkite, burna, laikas, baras.} \\
\texttt{Jūsų priešininko komandos žodžiai yra: styga, pyragas, amerika, mėta, siena, siela, degtukas, blokas.} \\
\texttt{Atitraukimo žodžiai yra: jupiteris, erelis, šepetys.} \\
\texttt{Žudymo žodžiai yra: alyvuogė, lokys, kengūra, slaugytoja, maskva.} \\
            }
        }
    }
    & \\ \\

    \theutterance \stepcounter{utterance}  
    & \multicolumn{4}{p{0.6\linewidth}}{
        \cellcolor[rgb]{0.9,0.9,0.9}{
            \makecell[{{p{\linewidth}}}]{
                \texttt{\tiny{[P1$\rangle$GM]}}
                \texttt{UŽUOMINA: saldus} \\
\texttt{TIKSLAI: medus, morka, baras} \\
            }
        }
    }
    & & \\ \\

    \theutterance \stepcounter{utterance}  
    & & & \multicolumn{2}{p{0.3\linewidth}}{
        \cellcolor[rgb]{0.9,0.9,0.9}{
            \makecell[{{p{\linewidth}}}]{
                \texttt{\tiny{[GM$|$GM]}}
                \texttt{saldus} \\
            }
        }
    }
    & & \\ \\

    \theutterance \stepcounter{utterance}  
    & & & \multicolumn{2}{p{0.3\linewidth}}{
        \cellcolor[rgb]{0.9,0.9,0.9}{
            \makecell[{{p{\linewidth}}}]{
                \texttt{\tiny{[GM$|$GM]}}
                ['medus', 'morka', 'baras']
            }
        }
    }
    & & \\ \\

    \theutterance \stepcounter{utterance}  
    & & & \multicolumn{2}{p{0.3\linewidth}}{
        \cellcolor[rgb]{0.9,0.9,0.9}{
            \makecell[{{p{\linewidth}}}]{
                \texttt{\tiny{[GM$|$GM]}}
                {'word': 'medus', 'assignment': 'team'}
            }
        }
    }
    & & \\ \\

    \theutterance \stepcounter{utterance}  
    & & & \multicolumn{2}{p{0.3\linewidth}}{
        \cellcolor[rgb]{0.9,0.9,0.9}{
            \makecell[{{p{\linewidth}}}]{
                \texttt{\tiny{[GM$|$GM]}}
                {'word': 'morka', 'assignment': 'team'}
            }
        }
    }
    & & \\ \\

    \theutterance \stepcounter{utterance}  
    & & & \multicolumn{2}{p{0.3\linewidth}}{
        \cellcolor[rgb]{0.9,0.9,0.9}{
            \makecell[{{p{\linewidth}}}]{
                \texttt{\tiny{[GM$|$GM]}}
                {'word': 'baras', 'assignment': 'team'}
            }
        }
    }
    & & \\ \\

    \theutterance \stepcounter{utterance}  
    & & \multicolumn{4}{p{0.6\linewidth}}{
        \cellcolor[rgb]{0.9,0.9,0.9}{
            \makecell[{{p{\linewidth}}}]{
                \texttt{\tiny{[GM$\rangle$P2]}}
                \texttt{Pateikite kableliais atskirtą sąrašą iki 3 žodžių iš toliau pateikto sąrašo, kurie geriausiai siejasi arba yra labiausiai susiję su žodžiu 'saldus'. Visada pradėkite savo spėjimų sąrašą žodžiais SPĖJIMAS:  ir į atsakymą neįtraukite jokio kito teksto.} \\
\\ 
\texttt{jupiteris, styga, pyragas, siela, maskva, erelis, blokas, kasykla, siena, alyvuogė, lochnesas, medus, lokys, mėta, degtukas, bakstelėkite, šepetys, arklys, burna, slaugytoja, amerika, morka, baras, laikas, kengūra} \\
            }
        }
    }
    & \\ \\

    \theutterance \stepcounter{utterance}  
    & & & \multicolumn{4}{p{0.6\linewidth}}{
        \cellcolor[rgb]{0.9,0.9,0.9}{
            \makecell[{{p{\linewidth}}}]{
                \texttt{\tiny{[GM$\langle$P2]}}
                \texttt{SPĖJIMAS: pyragas, medus, mėta} \\
            }
        }
    }
     \\ \\

    \theutterance \stepcounter{utterance}  
    & & & \multicolumn{2}{p{0.3\linewidth}}{
        \cellcolor[rgb]{0.9,0.9,0.9}{
            \makecell[{{p{\linewidth}}}]{
                \texttt{\tiny{[GM$|$GM]}}
                ['pyragas', 'medus', 'mėta']
            }
        }
    }
    & & \\ \\

    \theutterance \stepcounter{utterance}  
    & & & \multicolumn{2}{p{0.3\linewidth}}{
        \cellcolor[rgb]{0.9,0.9,0.9}{
            \makecell[{{p{\linewidth}}}]{
                \texttt{\tiny{[GM$|$GM]}}
                {'word': 'pyragas', 'assignment': 'opponent'}
            }
        }
    }
    & & \\ \\

    \theutterance \stepcounter{utterance}  
    & & & \multicolumn{2}{p{0.3\linewidth}}{
        \cellcolor[rgb]{0.9,0.9,0.9}{
            \makecell[{{p{\linewidth}}}]{
                \texttt{\tiny{[GM$|$GM]}}
                \texttt{pyragas} \\
            }
        }
    }
    & & \\ \\

    \theutterance \stepcounter{utterance}  
    & & & \multicolumn{2}{p{0.3\linewidth}}{
        \cellcolor[rgb]{0.9,0.9,0.9}{
            \makecell[{{p{\linewidth}}}]{
                \texttt{\tiny{[GM$|$GM]}}
                {'word': 'blokas', 'assignment': 'opponent'}
            }
        }
    }
    & & \\ \\

    \theutterance \stepcounter{utterance}  
    & & & \multicolumn{2}{p{0.3\linewidth}}{
        \cellcolor[rgb]{0.9,0.9,0.9}{
            \makecell[{{p{\linewidth}}}]{
                \texttt{\tiny{[GM$|$GM]}}
                {'hidden': {'team': ['medus', 'lochnesas', 'kasykla', 'arklys', 'morka', 'bakstelėkite', 'burna', 'laikas', 'baras'], 'innocent': ['jupiteris', 'erelis', 'šepetys'], 'opponent': ['styga', 'amerika', 'mėta', 'siena', 'siela', 'degtukas'], 'assassin': ['alyvuogė', 'lokys', 'kengūra', 'slaugytoja', 'maskva']}, 'revealed': {'team': {'team': [], 'innocent': [], 'opponent': ['pyragas'], 'assassin': []}, 'opponent': {'team': [], 'innocent': [], 'opponent': ['blokas'], 'assassin': []}}}
            }
        }
    }
    & & \\ \\

    \theutterance \stepcounter{utterance}  
    & & \multicolumn{4}{p{0.6\linewidth}}{
        \cellcolor[rgb]{0.9,0.9,0.9}{
            \makecell[{{p{\linewidth}}}]{
                \texttt{\tiny{[P1$\langle$GM]}}
                \texttt{Žodis pyragas buvo atspėtas, tačiau tai priešas žodis. Jūsų komandos draugo ėjimas baigėsi.} \\
\texttt{Dabar pateikite kitą užuominą, susijusią su kai kuriais likusiais jūsų komandos žodžiais, ir susijusių tikslinių žodžių sąrašą. Nepamirškite užuominą pradėti žodžiais "CLUE: ", įrašykite naują eilutę, o tikslinių žodžių sąrašą, atskirtą kableliais, pradėkite žodžiais "TARGETS: ". Pastaba: palyginti su ankstesniais prašymais, kai kurie žodžiai iš sąrašų buvo pašalinti.} \\
\\ 
\texttt{Likę jūsų komandos žodžiai yra šie: medus, lochnesas, kasykla, arklys, morka, bakstelėkite, burna, laikas, baras.} \\
\texttt{Likę priešininko žodžiai yra šie: styga, amerika, mėta, siena, siela, degtukas, blokas.} \\
\texttt{Likę atitraukiantys žodžiai yra: jupiteris, erelis, šepetys.} \\
\texttt{Likę nužudymo žodžiai yra: alyvuogė, lokys, kengūra, slaugytoja, maskva.} \\
            }
        }
    }
    & \\ \\

    \theutterance \stepcounter{utterance}  
    & \multicolumn{4}{p{0.6\linewidth}}{
        \cellcolor[rgb]{0.9,0.9,0.9}{
            \makecell[{{p{\linewidth}}}]{
                \texttt{\tiny{[P1$\rangle$GM]}}
                \texttt{UŽUOMINA: geltonas} \\
\texttt{TIKSLAI: morka, arklys} \\
            }
        }
    }
    & & \\ \\

    \theutterance \stepcounter{utterance}  
    & & & \multicolumn{2}{p{0.3\linewidth}}{
        \cellcolor[rgb]{0.9,0.9,0.9}{
            \makecell[{{p{\linewidth}}}]{
                \texttt{\tiny{[GM$|$GM]}}
                \texttt{geltonas} \\
            }
        }
    }
    & & \\ \\

    \theutterance \stepcounter{utterance}  
    & & & \multicolumn{2}{p{0.3\linewidth}}{
        \cellcolor[rgb]{0.9,0.9,0.9}{
            \makecell[{{p{\linewidth}}}]{
                \texttt{\tiny{[GM$|$GM]}}
                ['morka', 'arklys']
            }
        }
    }
    & & \\ \\

    \theutterance \stepcounter{utterance}  
    & & & \multicolumn{2}{p{0.3\linewidth}}{
        \cellcolor[rgb]{0.9,0.9,0.9}{
            \makecell[{{p{\linewidth}}}]{
                \texttt{\tiny{[GM$|$GM]}}
                {'word': 'morka', 'assignment': 'team'}
            }
        }
    }
    & & \\ \\

    \theutterance \stepcounter{utterance}  
    & & & \multicolumn{2}{p{0.3\linewidth}}{
        \cellcolor[rgb]{0.9,0.9,0.9}{
            \makecell[{{p{\linewidth}}}]{
                \texttt{\tiny{[GM$|$GM]}}
                {'word': 'arklys', 'assignment': 'team'}
            }
        }
    }
    & & \\ \\

    \theutterance \stepcounter{utterance}  
    & & \multicolumn{4}{p{0.6\linewidth}}{
        \cellcolor[rgb]{0.9,0.9,0.9}{
            \makecell[{{p{\linewidth}}}]{
                \texttt{\tiny{[GM$\rangle$P2]}}
                \texttt{Žodis pyragas buvo atspėtas, tačiau tai priešas žodis. Jūsų ėjimas baigėsi.} \\
\texttt{Dabar pateikite dar vieną kableliais atskirtą sąrašą iš bent 1 ir ne daugiau kaip 2 žodžių iš toliau pateikto žodžių sąrašo, kurie geriausiai siejasi arba yra labiausiai susiję su žodžiu 'geltonas'. Nepamirškite savo atsakymą pradėti SPĖJIMAS: . Pastaba: palyginti su ankstesniais prašymais, kai kurie žodžiai iš sąrašo buvo pašalinti.} \\
\\ 
\texttt{jupiteris, styga, siela, maskva, erelis, kasykla, siena, alyvuogė, lochnesas, medus, lokys, mėta, degtukas, bakstelėkite, šepetys, arklys, burna, slaugytoja, amerika, morka, baras, laikas, kengūra} \\
            }
        }
    }
    & \\ \\

    \theutterance \stepcounter{utterance}  
    & & & \multicolumn{4}{p{0.6\linewidth}}{
        \cellcolor[rgb]{0.9,0.9,0.9}{
            \makecell[{{p{\linewidth}}}]{
                \texttt{\tiny{[GM$\langle$P2]}}
                \texttt{SPĖJIMAS: jupiteris, morka} \\
            }
        }
    }
     \\ \\

    \theutterance \stepcounter{utterance}  
    & & & \multicolumn{2}{p{0.3\linewidth}}{
        \cellcolor[rgb]{0.9,0.9,0.9}{
            \makecell[{{p{\linewidth}}}]{
                \texttt{\tiny{[GM$|$GM]}}
                ['jupiteris', 'morka']
            }
        }
    }
    & & \\ \\

    \theutterance \stepcounter{utterance}  
    & & & \multicolumn{2}{p{0.3\linewidth}}{
        \cellcolor[rgb]{0.9,0.9,0.9}{
            \makecell[{{p{\linewidth}}}]{
                \texttt{\tiny{[GM$|$GM]}}
                {'word': 'jupiteris', 'assignment': 'innocent'}
            }
        }
    }
    & & \\ \\

    \theutterance \stepcounter{utterance}  
    & & & \multicolumn{2}{p{0.3\linewidth}}{
        \cellcolor[rgb]{0.9,0.9,0.9}{
            \makecell[{{p{\linewidth}}}]{
                \texttt{\tiny{[GM$|$GM]}}
                \texttt{jupiteris} \\
            }
        }
    }
    & & \\ \\

    \theutterance \stepcounter{utterance}  
    & & & \multicolumn{2}{p{0.3\linewidth}}{
        \cellcolor[rgb]{0.9,0.9,0.9}{
            \makecell[{{p{\linewidth}}}]{
                \texttt{\tiny{[GM$|$GM]}}
                {'word': 'mėta', 'assignment': 'opponent'}
            }
        }
    }
    & & \\ \\

    \theutterance \stepcounter{utterance}  
    & & & \multicolumn{2}{p{0.3\linewidth}}{
        \cellcolor[rgb]{0.9,0.9,0.9}{
            \makecell[{{p{\linewidth}}}]{
                \texttt{\tiny{[GM$|$GM]}}
                {'hidden': {'team': ['medus', 'lochnesas', 'kasykla', 'arklys', 'morka', 'bakstelėkite', 'burna', 'laikas', 'baras'], 'innocent': ['erelis', 'šepetys'], 'opponent': ['styga', 'amerika', 'siena', 'siela', 'degtukas'], 'assassin': ['alyvuogė', 'lokys', 'kengūra', 'slaugytoja', 'maskva']}, 'revealed': {'team': {'team': [], 'innocent': ['jupiteris'], 'opponent': ['pyragas'], 'assassin': []}, 'opponent': {'team': [], 'innocent': [], 'opponent': ['blokas', 'mėta'], 'assassin': []}}}
            }
        }
    }
    & & \\ \\

    \theutterance \stepcounter{utterance}  
    & & \multicolumn{4}{p{0.6\linewidth}}{
        \cellcolor[rgb]{0.9,0.9,0.9}{
            \makecell[{{p{\linewidth}}}]{
                \texttt{\tiny{[P1$\langle$GM]}}
                \texttt{Žodis jupiteris buvo atspėtas, tačiau tai nekaltas žodis. Jūsų komandos draugo ėjimas baigėsi.} \\
\texttt{Dabar pateikite kitą užuominą, susijusią su kai kuriais likusiais jūsų komandos žodžiais, ir susijusių tikslinių žodžių sąrašą. Nepamirškite užuominą pradėti žodžiais "CLUE: ", įrašykite naują eilutę, o tikslinių žodžių sąrašą, atskirtą kableliais, pradėkite žodžiais "TARGETS: ". Pastaba: palyginti su ankstesniais prašymais, kai kurie žodžiai iš sąrašų buvo pašalinti.} \\
\\ 
\texttt{Likę jūsų komandos žodžiai yra šie: medus, lochnesas, kasykla, arklys, morka, bakstelėkite, burna, laikas, baras.} \\
\texttt{Likę priešininko žodžiai yra šie: styga, amerika, mėta, siena, siela, degtukas.} \\
\texttt{Likę atitraukiantys žodžiai yra: erelis, šepetys.} \\
\texttt{Likę nužudymo žodžiai yra: alyvuogė, lokys, kengūra, slaugytoja, maskva.} \\
            }
        }
    }
    & \\ \\

    \theutterance \stepcounter{utterance}  
    & \multicolumn{4}{p{0.6\linewidth}}{
        \cellcolor[rgb]{0.9,0.9,0.9}{
            \makecell[{{p{\linewidth}}}]{
                \texttt{\tiny{[P1$\rangle$GM]}}
                \texttt{UŽUOMINA: geltonas} \\
\texttt{TIKSLAI: morka, arklys} \\
            }
        }
    }
    & & \\ \\

    \theutterance \stepcounter{utterance}  
    & & & \multicolumn{2}{p{0.3\linewidth}}{
        \cellcolor[rgb]{0.9,0.9,0.9}{
            \makecell[{{p{\linewidth}}}]{
                \texttt{\tiny{[GM$|$GM]}}
                \texttt{geltonas} \\
            }
        }
    }
    & & \\ \\

    \theutterance \stepcounter{utterance}  
    & & & \multicolumn{2}{p{0.3\linewidth}}{
        \cellcolor[rgb]{0.9,0.9,0.9}{
            \makecell[{{p{\linewidth}}}]{
                \texttt{\tiny{[GM$|$GM]}}
                ['morka', 'arklys']
            }
        }
    }
    & & \\ \\

    \theutterance \stepcounter{utterance}  
    & & & \multicolumn{2}{p{0.3\linewidth}}{
        \cellcolor[rgb]{0.9,0.9,0.9}{
            \makecell[{{p{\linewidth}}}]{
                \texttt{\tiny{[GM$|$GM]}}
                {'word': 'morka', 'assignment': 'team'}
            }
        }
    }
    & & \\ \\

    \theutterance \stepcounter{utterance}  
    & & & \multicolumn{2}{p{0.3\linewidth}}{
        \cellcolor[rgb]{0.9,0.9,0.9}{
            \makecell[{{p{\linewidth}}}]{
                \texttt{\tiny{[GM$|$GM]}}
                {'word': 'arklys', 'assignment': 'team'}
            }
        }
    }
    & & \\ \\

    \theutterance \stepcounter{utterance}  
    & & \multicolumn{4}{p{0.6\linewidth}}{
        \cellcolor[rgb]{0.9,0.9,0.9}{
            \makecell[{{p{\linewidth}}}]{
                \texttt{\tiny{[GM$\rangle$P2]}}
                \texttt{Žodis jupiteris buvo atspėtas, tačiau tai nekaltas žodis. Jūsų ėjimas baigėsi.} \\
\texttt{Dabar pateikite dar vieną kableliais atskirtą sąrašą iš bent 1 ir ne daugiau kaip 2 žodžių iš toliau pateikto žodžių sąrašo, kurie geriausiai siejasi arba yra labiausiai susiję su žodžiu 'geltonas'. Nepamirškite savo atsakymą pradėti SPĖJIMAS: . Pastaba: palyginti su ankstesniais prašymais, kai kurie žodžiai iš sąrašo buvo pašalinti.} \\
\\ 
\texttt{styga, siela, maskva, erelis, kasykla, siena, alyvuogė, lochnesas, medus, lokys, degtukas, bakstelėkite, šepetys, arklys, burna, slaugytoja, amerika, morka, baras, laikas, kengūra} \\
            }
        }
    }
    & \\ \\

    \theutterance \stepcounter{utterance}  
    & & & \multicolumn{4}{p{0.6\linewidth}}{
        \cellcolor[rgb]{0.9,0.9,0.9}{
            \makecell[{{p{\linewidth}}}]{
                \texttt{\tiny{[GM$\langle$P2]}}
                \texttt{SPĖJIMAS: medus, morka} \\
            }
        }
    }
     \\ \\

    \theutterance \stepcounter{utterance}  
    & & & \multicolumn{2}{p{0.3\linewidth}}{
        \cellcolor[rgb]{0.9,0.9,0.9}{
            \makecell[{{p{\linewidth}}}]{
                \texttt{\tiny{[GM$|$GM]}}
                ['medus', 'morka']
            }
        }
    }
    & & \\ \\

    \theutterance \stepcounter{utterance}  
    & & & \multicolumn{2}{p{0.3\linewidth}}{
        \cellcolor[rgb]{0.9,0.9,0.9}{
            \makecell[{{p{\linewidth}}}]{
                \texttt{\tiny{[GM$|$GM]}}
                {'word': 'medus', 'assignment': 'team'}
            }
        }
    }
    & & \\ \\

    \theutterance \stepcounter{utterance}  
    & & & \multicolumn{2}{p{0.3\linewidth}}{
        \cellcolor[rgb]{0.9,0.9,0.9}{
            \makecell[{{p{\linewidth}}}]{
                \texttt{\tiny{[GM$|$GM]}}
                {'word': 'morka', 'assignment': 'team'}
            }
        }
    }
    & & \\ \\

    \theutterance \stepcounter{utterance}  
    & & & \multicolumn{2}{p{0.3\linewidth}}{
        \cellcolor[rgb]{0.9,0.9,0.9}{
            \makecell[{{p{\linewidth}}}]{
                \texttt{\tiny{[GM$|$GM]}}
                {'word': 'morka', 'assignment': 'team'}
            }
        }
    }
    & & \\ \\

    \theutterance \stepcounter{utterance}  
    & & & \multicolumn{2}{p{0.3\linewidth}}{
        \cellcolor[rgb]{0.9,0.9,0.9}{
            \makecell[{{p{\linewidth}}}]{
                \texttt{\tiny{[GM$|$GM]}}
                {'word': 'styga', 'assignment': 'opponent'}
            }
        }
    }
    & & \\ \\

    \theutterance \stepcounter{utterance}  
    & & & \multicolumn{2}{p{0.3\linewidth}}{
        \cellcolor[rgb]{0.9,0.9,0.9}{
            \makecell[{{p{\linewidth}}}]{
                \texttt{\tiny{[GM$|$GM]}}
                {'hidden': {'team': ['lochnesas', 'kasykla', 'arklys', 'bakstelėkite', 'burna', 'laikas', 'baras'], 'innocent': ['erelis', 'šepetys'], 'opponent': ['amerika', 'siena', 'siela', 'degtukas'], 'assassin': ['alyvuogė', 'lokys', 'kengūra', 'slaugytoja', 'maskva']}, 'revealed': {'team': {'team': ['medus', 'morka'], 'innocent': ['jupiteris'], 'opponent': ['pyragas'], 'assassin': []}, 'opponent': {'team': [], 'innocent': [], 'opponent': ['blokas', 'mėta', 'styga'], 'assassin': []}}}
            }
        }
    }
    & & \\ \\

    \theutterance \stepcounter{utterance}  
    & & \multicolumn{4}{p{0.6\linewidth}}{
        \cellcolor[rgb]{0.9,0.9,0.9}{
            \makecell[{{p{\linewidth}}}]{
                \texttt{\tiny{[P1$\langle$GM]}}
                \texttt{Žodžiai medus, morka buvo atspėti teisingai. Jūsų komandos draugo ėjimas baigėsi.} \\
\texttt{Dabar pateikite kitą užuominą, susijusią su kai kuriais likusiais jūsų komandos žodžiais, ir susijusių tikslinių žodžių sąrašą. Nepamirškite užuominą pradėti žodžiais "CLUE: ", įrašykite naują eilutę, o tikslinių žodžių sąrašą, atskirtą kableliais, pradėkite žodžiais "TARGETS: ". Pastaba: palyginti su ankstesniais prašymais, kai kurie žodžiai iš sąrašų buvo pašalinti.} \\
\\ 
\texttt{Likę jūsų komandos žodžiai yra šie: lochnesas, kasykla, arklys, bakstelėkite, burna, laikas, baras.} \\
\texttt{Likę priešininko žodžiai yra šie: styga, amerika, siena, siela, degtukas.} \\
\texttt{Likę atitraukiantys žodžiai yra: erelis, šepetys.} \\
\texttt{Likę nužudymo žodžiai yra: alyvuogė, lokys, kengūra, slaugytoja, maskva.} \\
            }
        }
    }
    & \\ \\

    \theutterance \stepcounter{utterance}  
    & \multicolumn{4}{p{0.6\linewidth}}{
        \cellcolor[rgb]{0.9,0.9,0.9}{
            \makecell[{{p{\linewidth}}}]{
                \texttt{\tiny{[P1$\rangle$GM]}}
                \texttt{UŽUOMINA: vanduo} \\
\texttt{TIKSLAI: lochnesas, kasykla} \\
            }
        }
    }
    & & \\ \\

    \theutterance \stepcounter{utterance}  
    & & & \multicolumn{2}{p{0.3\linewidth}}{
        \cellcolor[rgb]{0.9,0.9,0.9}{
            \makecell[{{p{\linewidth}}}]{
                \texttt{\tiny{[GM$|$GM]}}
                \texttt{vanduo} \\
            }
        }
    }
    & & \\ \\

    \theutterance \stepcounter{utterance}  
    & & & \multicolumn{2}{p{0.3\linewidth}}{
        \cellcolor[rgb]{0.9,0.9,0.9}{
            \makecell[{{p{\linewidth}}}]{
                \texttt{\tiny{[GM$|$GM]}}
                ['lochnesas', 'kasykla']
            }
        }
    }
    & & \\ \\

    \theutterance \stepcounter{utterance}  
    & & & \multicolumn{2}{p{0.3\linewidth}}{
        \cellcolor[rgb]{0.9,0.9,0.9}{
            \makecell[{{p{\linewidth}}}]{
                \texttt{\tiny{[GM$|$GM]}}
                {'word': 'lochnesas', 'assignment': 'team'}
            }
        }
    }
    & & \\ \\

    \theutterance \stepcounter{utterance}  
    & & & \multicolumn{2}{p{0.3\linewidth}}{
        \cellcolor[rgb]{0.9,0.9,0.9}{
            \makecell[{{p{\linewidth}}}]{
                \texttt{\tiny{[GM$|$GM]}}
                {'word': 'kasykla', 'assignment': 'team'}
            }
        }
    }
    & & \\ \\

    \theutterance \stepcounter{utterance}  
    & & \multicolumn{4}{p{0.6\linewidth}}{
        \cellcolor[rgb]{0.9,0.9,0.9}{
            \makecell[{{p{\linewidth}}}]{
                \texttt{\tiny{[GM$\rangle$P2]}}
                \texttt{Žodžiai medus, morka buvo atspėti teisingai. Jūsų ėjimas baigėsi.} \\
\texttt{Dabar pateikite dar vieną kableliais atskirtą sąrašą iš bent 1 ir ne daugiau kaip 2 žodžių iš toliau pateikto žodžių sąrašo, kurie geriausiai siejasi arba yra labiausiai susiję su žodžiu 'vanduo'. Nepamirškite savo atsakymą pradėti SPĖJIMAS: . Pastaba: palyginti su ankstesniais prašymais, kai kurie žodžiai iš sąrašo buvo pašalinti.} \\
\\ 
\texttt{siela, maskva, erelis, kasykla, siena, alyvuogė, lochnesas, lokys, degtukas, bakstelėkite, šepetys, arklys, burna, slaugytoja, amerika, baras, laikas, kengūra} \\
            }
        }
    }
    & \\ \\

    \theutterance \stepcounter{utterance}  
    & & & \multicolumn{4}{p{0.6\linewidth}}{
        \cellcolor[rgb]{0.9,0.9,0.9}{
            \makecell[{{p{\linewidth}}}]{
                \texttt{\tiny{[GM$\langle$P2]}}
                \texttt{SPĖJIMAS: siela, maskva} \\
            }
        }
    }
     \\ \\

    \theutterance \stepcounter{utterance}  
    & & & \multicolumn{2}{p{0.3\linewidth}}{
        \cellcolor[rgb]{0.9,0.9,0.9}{
            \makecell[{{p{\linewidth}}}]{
                \texttt{\tiny{[GM$|$GM]}}
                ['siela', 'maskva']
            }
        }
    }
    & & \\ \\

    \theutterance \stepcounter{utterance}  
    & & & \multicolumn{2}{p{0.3\linewidth}}{
        \cellcolor[rgb]{0.9,0.9,0.9}{
            \makecell[{{p{\linewidth}}}]{
                \texttt{\tiny{[GM$|$GM]}}
                {'word': 'siela', 'assignment': 'opponent'}
            }
        }
    }
    & & \\ \\

    \theutterance \stepcounter{utterance}  
    & & & \multicolumn{2}{p{0.3\linewidth}}{
        \cellcolor[rgb]{0.9,0.9,0.9}{
            \makecell[{{p{\linewidth}}}]{
                \texttt{\tiny{[GM$|$GM]}}
                \texttt{siela} \\
            }
        }
    }
    & & \\ \\

    \theutterance \stepcounter{utterance}  
    & & & \multicolumn{2}{p{0.3\linewidth}}{
        \cellcolor[rgb]{0.9,0.9,0.9}{
            \makecell[{{p{\linewidth}}}]{
                \texttt{\tiny{[GM$|$GM]}}
                {'word': 'amerika', 'assignment': 'opponent'}
            }
        }
    }
    & & \\ \\

    \theutterance \stepcounter{utterance}  
    & & & \multicolumn{2}{p{0.3\linewidth}}{
        \cellcolor[rgb]{0.9,0.9,0.9}{
            \makecell[{{p{\linewidth}}}]{
                \texttt{\tiny{[GM$|$GM]}}
                {'hidden': {'team': ['lochnesas', 'kasykla', 'arklys', 'bakstelėkite', 'burna', 'laikas', 'baras'], 'innocent': ['erelis', 'šepetys'], 'opponent': ['siena', 'degtukas'], 'assassin': ['alyvuogė', 'lokys', 'kengūra', 'slaugytoja', 'maskva']}, 'revealed': {'team': {'team': ['medus', 'morka'], 'innocent': ['jupiteris'], 'opponent': ['pyragas', 'siela'], 'assassin': []}, 'opponent': {'team': [], 'innocent': [], 'opponent': ['blokas', 'mėta', 'styga', 'amerika'], 'assassin': []}}}
            }
        }
    }
    & & \\ \\

    \theutterance \stepcounter{utterance}  
    & & \multicolumn{4}{p{0.6\linewidth}}{
        \cellcolor[rgb]{0.9,0.9,0.9}{
            \makecell[{{p{\linewidth}}}]{
                \texttt{\tiny{[P1$\langle$GM]}}
                \texttt{Žodis siela buvo atspėtas, tačiau tai priešas žodis. Jūsų komandos draugo ėjimas baigėsi.} \\
\texttt{Dabar pateikite kitą užuominą, susijusią su kai kuriais likusiais jūsų komandos žodžiais, ir susijusių tikslinių žodžių sąrašą. Nepamirškite užuominą pradėti žodžiais "CLUE: ", įrašykite naują eilutę, o tikslinių žodžių sąrašą, atskirtą kableliais, pradėkite žodžiais "TARGETS: ". Pastaba: palyginti su ankstesniais prašymais, kai kurie žodžiai iš sąrašų buvo pašalinti.} \\
\\ 
\texttt{Likę jūsų komandos žodžiai yra šie: lochnesas, kasykla, arklys, bakstelėkite, burna, laikas, baras.} \\
\texttt{Likę priešininko žodžiai yra šie: amerika, siena, degtukas.} \\
\texttt{Likę atitraukiantys žodžiai yra: erelis, šepetys.} \\
\texttt{Likę nužudymo žodžiai yra: alyvuogė, lokys, kengūra, slaugytoja, maskva.} \\
            }
        }
    }
    & \\ \\

    \theutterance \stepcounter{utterance}  
    & \multicolumn{4}{p{0.6\linewidth}}{
        \cellcolor[rgb]{0.9,0.9,0.9}{
            \makecell[{{p{\linewidth}}}]{
                \texttt{\tiny{[P1$\rangle$GM]}}
                \texttt{UŽUOMINA: gyvūnas} \\
\texttt{TIKSLAI: arklys} \\
            }
        }
    }
    & & \\ \\

    \theutterance \stepcounter{utterance}  
    & & & \multicolumn{2}{p{0.3\linewidth}}{
        \cellcolor[rgb]{0.9,0.9,0.9}{
            \makecell[{{p{\linewidth}}}]{
                \texttt{\tiny{[GM$|$GM]}}
                \texttt{gyvūnas} \\
            }
        }
    }
    & & \\ \\

    \theutterance \stepcounter{utterance}  
    & & & \multicolumn{2}{p{0.3\linewidth}}{
        \cellcolor[rgb]{0.9,0.9,0.9}{
            \makecell[{{p{\linewidth}}}]{
                \texttt{\tiny{[GM$|$GM]}}
                ['arklys']
            }
        }
    }
    & & \\ \\

    \theutterance \stepcounter{utterance}  
    & & & \multicolumn{2}{p{0.3\linewidth}}{
        \cellcolor[rgb]{0.9,0.9,0.9}{
            \makecell[{{p{\linewidth}}}]{
                \texttt{\tiny{[GM$|$GM]}}
                {'word': 'arklys', 'assignment': 'team'}
            }
        }
    }
    & & \\ \\

    \theutterance \stepcounter{utterance}  
    & & \multicolumn{4}{p{0.6\linewidth}}{
        \cellcolor[rgb]{0.9,0.9,0.9}{
            \makecell[{{p{\linewidth}}}]{
                \texttt{\tiny{[GM$\rangle$P2]}}
                \texttt{Žodis siela buvo atspėtas, tačiau tai priešas žodis. Jūsų ėjimas baigėsi.} \\
\texttt{Dabar pateikite dar vieną kableliais atskirtą sąrašą iš bent 1 ir ne daugiau kaip 1 žodžių iš toliau pateikto žodžių sąrašo, kurie geriausiai siejasi arba yra labiausiai susiję su žodžiu 'gyvūnas'. Nepamirškite savo atsakymą pradėti SPĖJIMAS: . Pastaba: palyginti su ankstesniais prašymais, kai kurie žodžiai iš sąrašo buvo pašalinti.} \\
\\ 
\texttt{maskva, erelis, kasykla, siena, alyvuogė, lochnesas, lokys, degtukas, bakstelėkite, šepetys, arklys, burna, slaugytoja, baras, laikas, kengūra} \\
            }
        }
    }
    & \\ \\

    \theutterance \stepcounter{utterance}  
    & & & \multicolumn{4}{p{0.6\linewidth}}{
        \cellcolor[rgb]{0.9,0.9,0.9}{
            \makecell[{{p{\linewidth}}}]{
                \texttt{\tiny{[GM$\langle$P2]}}
                \texttt{SPĖJIMAS: erelis} \\
            }
        }
    }
     \\ \\

    \theutterance \stepcounter{utterance}  
    & & & \multicolumn{2}{p{0.3\linewidth}}{
        \cellcolor[rgb]{0.9,0.9,0.9}{
            \makecell[{{p{\linewidth}}}]{
                \texttt{\tiny{[GM$|$GM]}}
                ['erelis']
            }
        }
    }
    & & \\ \\

    \theutterance \stepcounter{utterance}  
    & & & \multicolumn{2}{p{0.3\linewidth}}{
        \cellcolor[rgb]{0.9,0.9,0.9}{
            \makecell[{{p{\linewidth}}}]{
                \texttt{\tiny{[GM$|$GM]}}
                {'word': 'erelis', 'assignment': 'innocent'}
            }
        }
    }
    & & \\ \\

    \theutterance \stepcounter{utterance}  
    & & & \multicolumn{2}{p{0.3\linewidth}}{
        \cellcolor[rgb]{0.9,0.9,0.9}{
            \makecell[{{p{\linewidth}}}]{
                \texttt{\tiny{[GM$|$GM]}}
                \texttt{erelis} \\
            }
        }
    }
    & & \\ \\

    \theutterance \stepcounter{utterance}  
    & & & \multicolumn{2}{p{0.3\linewidth}}{
        \cellcolor[rgb]{0.9,0.9,0.9}{
            \makecell[{{p{\linewidth}}}]{
                \texttt{\tiny{[GM$|$GM]}}
                {'word': 'degtukas', 'assignment': 'opponent'}
            }
        }
    }
    & & \\ \\

    \theutterance \stepcounter{utterance}  
    & & & \multicolumn{2}{p{0.3\linewidth}}{
        \cellcolor[rgb]{0.9,0.9,0.9}{
            \makecell[{{p{\linewidth}}}]{
                \texttt{\tiny{[GM$|$GM]}}
                {'hidden': {'team': ['lochnesas', 'kasykla', 'arklys', 'bakstelėkite', 'burna', 'laikas', 'baras'], 'innocent': ['šepetys'], 'opponent': ['siena'], 'assassin': ['alyvuogė', 'lokys', 'kengūra', 'slaugytoja', 'maskva']}, 'revealed': {'team': {'team': ['medus', 'morka'], 'innocent': ['jupiteris', 'erelis'], 'opponent': ['pyragas', 'siela'], 'assassin': []}, 'opponent': {'team': [], 'innocent': [], 'opponent': ['blokas', 'mėta', 'styga', 'amerika', 'degtukas'], 'assassin': []}}}
            }
        }
    }
    & & \\ \\

    \theutterance \stepcounter{utterance}  
    & & \multicolumn{4}{p{0.6\linewidth}}{
        \cellcolor[rgb]{0.9,0.9,0.9}{
            \makecell[{{p{\linewidth}}}]{
                \texttt{\tiny{[P1$\langle$GM]}}
                \texttt{Žodis erelis buvo atspėtas, tačiau tai nekaltas žodis. Jūsų komandos draugo ėjimas baigėsi.} \\
\texttt{Dabar pateikite kitą užuominą, susijusią su kai kuriais likusiais jūsų komandos žodžiais, ir susijusių tikslinių žodžių sąrašą. Nepamirškite užuominą pradėti žodžiais "CLUE: ", įrašykite naują eilutę, o tikslinių žodžių sąrašą, atskirtą kableliais, pradėkite žodžiais "TARGETS: ". Pastaba: palyginti su ankstesniais prašymais, kai kurie žodžiai iš sąrašų buvo pašalinti.} \\
\\ 
\texttt{Likę jūsų komandos žodžiai yra šie: lochnesas, kasykla, arklys, bakstelėkite, burna, laikas, baras.} \\
\texttt{Likę priešininko žodžiai yra šie: siena, degtukas.} \\
\texttt{Likę atitraukiantys žodžiai yra: šepetys.} \\
\texttt{Likę nužudymo žodžiai yra: alyvuogė, lokys, kengūra, slaugytoja, maskva.} \\
            }
        }
    }
    & \\ \\

    \theutterance \stepcounter{utterance}  
    & \multicolumn{4}{p{0.6\linewidth}}{
        \cellcolor[rgb]{0.9,0.9,0.9}{
            \makecell[{{p{\linewidth}}}]{
                \texttt{\tiny{[P1$\rangle$GM]}}
                \texttt{UŽUOMINA: gyvūnas} \\
\texttt{TIKSLAI: arklys} \\
            }
        }
    }
    & & \\ \\

    \theutterance \stepcounter{utterance}  
    & & & \multicolumn{2}{p{0.3\linewidth}}{
        \cellcolor[rgb]{0.9,0.9,0.9}{
            \makecell[{{p{\linewidth}}}]{
                \texttt{\tiny{[GM$|$GM]}}
                \texttt{gyvūnas} \\
            }
        }
    }
    & & \\ \\

    \theutterance \stepcounter{utterance}  
    & & & \multicolumn{2}{p{0.3\linewidth}}{
        \cellcolor[rgb]{0.9,0.9,0.9}{
            \makecell[{{p{\linewidth}}}]{
                \texttt{\tiny{[GM$|$GM]}}
                ['arklys']
            }
        }
    }
    & & \\ \\

    \theutterance \stepcounter{utterance}  
    & & & \multicolumn{2}{p{0.3\linewidth}}{
        \cellcolor[rgb]{0.9,0.9,0.9}{
            \makecell[{{p{\linewidth}}}]{
                \texttt{\tiny{[GM$|$GM]}}
                {'word': 'arklys', 'assignment': 'team'}
            }
        }
    }
    & & \\ \\

    \theutterance \stepcounter{utterance}  
    & & \multicolumn{4}{p{0.6\linewidth}}{
        \cellcolor[rgb]{0.9,0.9,0.9}{
            \makecell[{{p{\linewidth}}}]{
                \texttt{\tiny{[GM$\rangle$P2]}}
                \texttt{Žodis erelis buvo atspėtas, tačiau tai nekaltas žodis. Jūsų ėjimas baigėsi.} \\
\texttt{Dabar pateikite dar vieną kableliais atskirtą sąrašą iš bent 1 ir ne daugiau kaip 1 žodžių iš toliau pateikto žodžių sąrašo, kurie geriausiai siejasi arba yra labiausiai susiję su žodžiu 'gyvūnas'. Nepamirškite savo atsakymą pradėti SPĖJIMAS: . Pastaba: palyginti su ankstesniais prašymais, kai kurie žodžiai iš sąrašo buvo pašalinti.} \\
\\ 
\texttt{maskva, kasykla, siena, alyvuogė, lochnesas, lokys, bakstelėkite, šepetys, arklys, burna, slaugytoja, baras, laikas, kengūra} \\
            }
        }
    }
    & \\ \\

    \theutterance \stepcounter{utterance}  
    & & & \multicolumn{4}{p{0.6\linewidth}}{
        \cellcolor[rgb]{0.9,0.9,0.9}{
            \makecell[{{p{\linewidth}}}]{
                \texttt{\tiny{[GM$\langle$P2]}}
                \texttt{SPĖJIMAS: lokys} \\
            }
        }
    }
     \\ \\

    \theutterance \stepcounter{utterance}  
    & & & \multicolumn{2}{p{0.3\linewidth}}{
        \cellcolor[rgb]{0.9,0.9,0.9}{
            \makecell[{{p{\linewidth}}}]{
                \texttt{\tiny{[GM$|$GM]}}
                ['lokys']
            }
        }
    }
    & & \\ \\

    \theutterance \stepcounter{utterance}  
    & & & \multicolumn{2}{p{0.3\linewidth}}{
        \cellcolor[rgb]{0.9,0.9,0.9}{
            \makecell[{{p{\linewidth}}}]{
                \texttt{\tiny{[GM$|$GM]}}
                {'word': 'lokys', 'assignment': 'assassin'}
            }
        }
    }
    & & \\ \\

    \theutterance \stepcounter{utterance}  
    & & & \multicolumn{2}{p{0.3\linewidth}}{
        \cellcolor[rgb]{0.9,0.9,0.9}{
            \makecell[{{p{\linewidth}}}]{
                \texttt{\tiny{[GM$|$GM]}}
                \texttt{lokys} \\
            }
        }
    }
    & & \\ \\

    \theutterance \stepcounter{utterance}  
    & & & \multicolumn{2}{p{0.3\linewidth}}{
        \cellcolor[rgb]{0.9,0.9,0.9}{
            \makecell[{{p{\linewidth}}}]{
                \texttt{\tiny{[GM$|$GM]}}
                \texttt{opponent has won through assassin} \\
            }
        }
    }
    & & \\ \\

\end{supertabular}
}

\end{document}
